\section{Анализ существующих программных решений}
В данном разделе представлен обзор 30 реализаций систем для автоматизированной проверке чертежей на ошибки оформления, приведены критерии их сравнения и построена сравнительная таблица для найденых решений.
 
\subsection{Плагины и функции встроенные в САПР}

Существуют системы автоматизированного нормоконтроля чертежей, интегрированные в среды САПР.
Они взаимодействуют с САПР с помощью API и получают данные о чертежах без искажений.
Анализ таких систем, как КОМПАС-3D Эксперт~\cite{compas-ecspert}, AutoCAD Drawing Checker~\cite{autocad-drawing-checker}, Autodesk Design Checker~\cite{autodesk-checker},
SolidWorks Design Checker~\cite{solidwork-checker}, PTC Creo ModelCHECK~\cite{creo-checker}, Sie\-mens NX Check-Mate~\cite{nx-checker}, CATIA Design Rule Checker~\cite{catia-checker}, DraftSight Standards Checker~\cite{drafsite-checker}, 
T-FLEX CAD Model Check~\cite{t-flex-checker}, nanoCAD Standards~\cite{nanocad-checker}, CADSTA Standards Manage~~\cite{cadsta-standards}, а также корпоративные решения, например АСКОН ЛОЦМАН:PLM~\cite{ascon-locman},
выявляет общую архитектурную и функциональную направленность.
% 12 источников

Общей чертой большинства систем является подход, ориентированный на соблюдение формальных правил,
при котором проверка осуществляется на основе набора предопределённых или настраиваемых правил,
охватывающих такие параметры, как:
\begin{itemize}
    \item структура и иерархия слоёв;
    \item оформление размерных надписей и допусков;
    \item соответствие шрифтов и стилей аннотаций;
    \item корректность основной надписи;
    \item оформления рамки;
    \item соответствие геометрии модели заданным ограничениям.
\end{itemize}

У данного вида систем можно выделить следующие приемущества:
\begin{itemize}
    \item высокую степень автоматизации, устраняющую необходимость ручного перебора чертежа на соответствие требованиям;
    \item интеграцию с CAD-средой, что обеспечивает оперативную обратную связь и минимизирует временные затраты на корректировку;
    \item возможность настройки пользовательских правил, что расширяет применимость в корпоративных средах.
\end{itemize}
Однако наблюдается и ряд недостатков:
\begin{itemize}
    \item ограниченная семантическая глубина анализа, большинство систем не способны интерпретировать смысловую корректность или логические противоречия в проекте;
    \item зависимость от точности и полноты заданных правил.
\end{itemize}

\subsubsection{Altium Designer}

Altium Designer~\cite{altium-designer} --- специализированная САПР для проектирования печатных плат,
включающая встроенный Design Rule Check.
Проверяются геометрические и электрические правила:
\begin{itemize}
    \item минимальные зазоры между дорожками;
    \item целостность цепей;
    \item соответствие контактных площадок требованиям монтажа;
    \item корректность обозначений компонентов и слоёв.
\end{itemize}

\subsection{Системы автоматизированной проверки BIM моделей}

Современные решения для верификации BIM-моделей реализуют подход ориентированный на правила при
контроле проектной документации.
Были выделены платформы: Autodesk Model Checker for Revit~\cite{autodesk-revit-model-checker}, Solibri~\cite{solibri}, BIMcollab Zoom~{bimcollab-zoom} и Bentley Systems Standarts~\cite{bentley-standarts} ---
обеспечивают автоматизированную проверку моделей на соответствие стандартам.
Все три системы поддерживают выполнение настраиваемых правил,
охватывающих такие аспекты, как наличие и корректность атрибутов объектов, соблюдение пространственной структуры, а также
проверка геометрических и семантических свойств, наличие коллизий и заполненность обязательных параметров.

Отличительные особенности решений:
\begin{itemize}
    \item Autodesk Model Checker for Revit тесно интегрирован в среду Revit, позволяет напрямую переходить к несоответствующим элементам и поддерживает экспорт результатов для последующего анализа в Power BI;
    \item Solibri предлагает расширенную библиотеку проверок, включая анализ свободного пространства, специализированные проверки, например эвакуационные пути, и возможность построения сложных правил;
    \item BIMcollab Zoom акцентирует внимание на управлении жизненным циклом выявленных замечаний.
\end{itemize}
% 14

\subsection{AI-системы}
\subsubsection{Dessia}
Dessia~\cite{dessia} --- это программное решение на основе искусственного интеллекта, разработанное для автоматизированной верификации полноты и согласованности технической документации в инженерных проектах. Система специализируется на сопоставлении содержимого 2D-чертежей с документацией проекта: ведомостями материалов, контрольными листами и метаданными основной надписи. Валидация осуществляется по следующим ключевым параметрам:
\begin{itemize}
    \item распознавание и извлечение меток компонентов на 2D-чертежах;
    \item сопоставление позиций компонентов с указанными в ведомости материалов зонами;
    \item проверка корректности основной надписи;
    \item cемантическая сверка между чертежом и спецификацией: соответствие количества, обозначений и расположения компонентов требованиям ведомости материалов.
\end{itemize} 

Система использует гибридный подход, объединяющий методы компьютерного зрения для анализа векторных и растровых чертежей и правила логической валидации для выявления несоответствий. Результаты проверки визуализируются с помощью цветовой маркировки и формируются в отчёт, что позволяет сократить время ручной верификации. 

\subsubsection{Siemens NX AI Validation} 

Интегрирован в NX Check-Mate; использует обучаемые модели для семантического анализа аннотаций и выявления несоответствий с корпоративными шаблонами.~\cite{nx-ai-validation}.

\subsubsection{ThingWorx Analytics}

ThingWorx Analytics~\cite{thing-worx} —-- это компонент платформы ThingWorx от PTC~\cite{creo-checker}, предназначенный для интеграции методов искусственного интеллекта и машинного обучения в промышленные решения. Он может применяться как часть расширенной системы контроля Creo ModelCHECK, где используется для выявления аномалий и паттернов ошибок в оформлении чертежей и 3D-моделей на основе исторических данных проектов. Система автоматически формирует предиктивные модели, позволяя прогнозировать типичные нарушения и рекомендовать их устранение. 

\subsubsection{AI-Norma}

Norma.AI~\cite{norma-ai} --- это ИИ-система для автоматизированной пред экспертной проверки проектной документации.
Система осуществляет многоуровневый анализ чертежей и текстовых разделов на соответствие нормативным документам, поддерживая форматы DWG, PDF и другие.
Для проверки чертежей существуют функции: контроль размеров, сопоставление графических данных с пояснительной запиской и BIM/CAD-верификация. 

\subsubsection{DrawCheck}

DrawCheck.ai~\cite{drawcheck-ai} использует искусственный интеллект для автоматического анализа технических чертежей на соответствие отраслевым стандартам и внутренним корпоративным требованиям.
Система генерирует детализированные отчёты с оценкой соответствия и конкретными рекомендациями по исправлению выявленных несоответствий.

% 22

\subsection{Академические и исследовательские прототипы}

\subsubsection{Метод контроля твердотельных электронных моделей машиностроительного издания}

Метод~\cite{blinova2012} автоматической проверки решений конструктивных задач инженерной геометрии,
предложенный Бойковым и Федотовым, основан на сопоставлении проверяемой модели с эталоном,
заданным в виде формальной грамматики.
Система принимает чертёж в формате DXF, интерпретирует его как конструктивно-геометрическую
модель, проводит синтаксический разбор по контекстно-свободной грамматике и вычисляет меру различия между решением и эталоном.
Метод поддерживает задачи с множеством корректных решений и ориентирован на интеграцию в системы дистанционного обучения.

\subsubsection{Внедрение средств автоматической проверки решений конструктивных задач инженерной геометрии в CAD-систему}

Метод реализует автоматическую проверку решений конструктивных задач инженерной геометрии непосредственно в среде CAD-системы (на примере «КОМПАС-3D»),
что устраняет потери геометрической информации и экономит время, освобождая от необходимости экспортировать файлы.
Программная реализация выполнена в виде внешнего модуля на языке Delphi, использующего API «КОМПАС-3D» для прямого доступа к геометрической модели.
Модуль формирует конструктивную геометрическую модель решения, передаёт её по сети удалённой проверяющей системе, а затем визуализирует результаты проверки в интерфейсе CAD-системы~\cite{integrations-cad}.

\subsubsection{Применение нейросетей для оптимизации нормоконтроля}

Публикация Пеньшина И.С.~\cite{ai-method} посвящена разработке концепции системы автоматизированного нормоконтроля
на основе нейросетевых технологий и компьютерного зрения для проверки конструкторской документации
на соответствие требованиям ГОСТ и ЕСКД.
Предлагаемая система интегрируется с САПР,
выполняет предобработку чертежей, их анализ с помощью свёрточных нейронных сетей,
сопоставление с эталонами из нормативных документов и автоматическую генерацию аннотированных отчётов о несоответствиях.

\subsubsection{Метод конструктивно-геометрической проверки наложением}

Метод конструктивно-геометрической проверки наложением~\cite{boikov2016} реализует автоматизированный контроль чертежей
путём сопоставления проверяемой модели с эталонной через пространственное совмещение их геометрических объектов.
Система выявляет совпадающие и отсутствующие элементы, обеспечивая количественную оценку корректности решения
на основе взвешенной суммы правильно построенных объектов.

\subsubsection{Автоматическое извлечение информации из титульных блоков}

Метод автоматического обнаружению и извлечения информации из титульных блоков строительных чертежей~\cite{duan2024automatic}
с помощью гибридного подхода: Faster R-CNN используется для определения титульного блока на изображении,
а GPT-4o --- для извлечения структурированного текста из обнаруженной области,
это эффективно для бумажных чертежей: исторических, сканированных или рукописных.

\subsubsection{Автоматическая проверка чертежей для одного класса деталей}

В методе~\cite{boikov2020} предложено искать ошибки в чертежах посредством грамматического анализа геометрической формы детали,
основанного на заранее заданной контекстно-свободной грамматике,
описывающей допустимую последовательность элементарных форм (цилиндров, конусов, призм и т.п) и их модификаторов (фасок, резьбы и пр.).
Перед анализом выполняется предварительная обработка чертежа: устраняются дубликаты и наложения линий, выявляются геометрические ошибки
и разделяются изображения (виды, разрезы, сечения).
Затем с помощью нисходящего грамматического разбора система последовательно сопоставляет фрагменты чертежа с правилами грамматики ---
если разбор не удаётся завершить или остаются нераспознанные объекты,
это свидетельствует об ошибке в структуре или оформлении чертежа.
Таким образом, ошибки выявляются не только как нарушения ЕСКД,
но и как несоответствие логики изображения принятой геометрической модели класса деталей.
%28

\subsubsection{Автоматизированная система нормоконтроля конструкторской документации}

Метод~\cite{frantsuzova2017automated} реализует экспертную систему, основанную на базе типовых конструктивных решений,
включающей как корректные примеры, так и антипаттерны.
Выбор решения осуществляется через дерево вывода в диалоговом режиме.
Система проверяет соответствие проектируемой конструкции требованиям стандартов и типовым практикам
на этапе разработки, что позволяет выявлять ошибки до технического контроля.

\subsubsection{Автоматическое обнаружение геометрических ошибок на машиностроительных 2D-чертежах}

В методе~\cite{vasinauto} предлагается метод автоматизированного выявления геометрических ошибок,
основанный на формализованном представлении чертежа как набора векторных геометрических примитивов
и применении алгоритмов их взаимного анализа.
Метод выделяет три основные категории ошибок: наложения или совпадения примитивов,
нарушения параллельности и ортогональности и нарушения топологии: разрывы контуров, отсутствие или избыток примыканий.
Для каждой категории разработаны специализированные алгоритмы, которые сравнивают параметры примитивов
с заданными допусками и помечают недопустимые конфигурации.
Подход ориентирован на работу с векторными форматами САПР.

\subsection{Сравнительный анализ методов}
Для проведения сравнительного анализа были выбраны следующие критерии:
\begin{itemize}
    \item \textbf{Глубина анализа} --- насколько система способна выявлять не только формальные ошибки, но и смысловые. От 1 до 3, где 1 это низко, а 3 это высоко;
    \item \textbf{Возможность настройки} --- возможно ли добавление собственных правил. От 1 до 3, где 1 --- невозможно, 3 --- возможность полной настройки.
    \item \textbf{Устойчивость к ошибкам} --- как система ведёт себя при испорченных данных: сканированных чертежах, неточных линиях и т.д. От 1 до 3, где 1 --- низкая, 3 --- высокая.
    \item \textbf{Интеграция с САПР} --- поддерживает ли система работу внутри САПР. от 1 до 3, где 1 --- внешний сервис, 2 --- плагин, 3 --- встроенно.
\end{itemize}

Таблица сравнения проанализированных систем представлена в таблице~\ref{tab:comparison_full}.

{
\small
\begin{longtable}{|p{3.85cm}|p{2.75cm}|p{2.75cm}|p{2.75cm}|p{2.75cm}|}
\caption{Сравнительная характеристика 30 систем автоматизированного анализа чертежей}
\label{tab:comparison_full} \\
\hline
\textbf{Система} & 
\textbf{Глубина анализа} & 
\textbf{Возможность настройки} &
\textbf{Устойчивость к ошибкам} &
\textbf{Интеграция} \\
\hline
\endfirsthead

\multicolumn{5}{c}{Продолжение таблицы \ref{tab:comparison_full}} \\
\hline
\textbf{Система} & 
\textbf{Глубина анализа} & 
\textbf{Возможность настройки} &
\textbf{Устойчивость к ошибкам} &
\textbf{Интеграция} \\
\hline
\endhead

\hline
\endlastfoot

% Содержимое таблицы
КОМПАС-3D Эксперт & 1 & 3 & 3 & 3 \\
\hline
AutoCAD Drawing Checker & 1 & 3 & 3 & 2 \\
\hline
Autodesk Design Checker & 1 & 3 & 3 & 2 \\
\hline
SolidWorks Design Checker & 1 & 3 & 3 & 3 \\
\hline
PTC Creo ModelCHECK & 1 & 3 & 3 & 3 \\
\hline
Siemens NX Check-Mate & 1 & 3 & 3 & 3 \\
\hline
CATIA Design Rule Checker & 1 & 3 & 3 & 3 \\
\hline
DraftSight Standards Checker & 1 & 3 & 3 & 2 \\
\hline
T-FLEX CAD Model Check & 1 & 3 & 3 & 3 \\
\hline
nanoCAD Standards & 1 & 3 & 3 & 2 \\
\hline
CADSTA Standards Manager & 1 & 3 & 3 & 1 \\
\hline
АСКОН ЛОЦМАН:PLM & 1 & 3 & 3 & 3 \\
\hline
Altium Designer & 1 & 3 & 3 & 3 \\
\hline
Autodesk Model Checker & 2 & 3 & 3 & 3 \\
\hline
Solibri & 2 & 3 & 3 & 2 \\
\hline
BIMcollab Zoom & 2 & 3 & 3 & 2 \\
\hline
Bentley Systems Standards & 2 & 2 & 3 & 3 \\
\hline
Dessia & 3 & 2 & 2 & 1 \\
\hline
Siemens NX AI Validation & 2 & 2 & 2 & 3 \\
\hline
ThingWorx Analytics & 3 & 2 & 2 & 2 \\
\hline
Norma.AI & 3 & 2 & 2 & 1 \\
\hline
DrawCheck.ai & 3 & 2 & 2 & 1 \\
\hline
Метод сравнения с эталоном & 2 & 1 & 2 & 1 \\
\hline
Внедрение автоматической проверки в систему CAD & 2 & 2 & 2 & 3 \\
\hline
Нейросети для оптимизации нормоконтроля & 3 & 2 & 2 & 1 \\
\hline
Метод проверки наложением & 2 & 1 & 2 & 1 \\
\hline
Анализ титульного блока & 1 & 2 & 1 & 1 \\
\hline
Автоматическая проверка чертежей деталей одного класса & 2 & 1 & 2 & 1 \\
\hline
Автоматизированная система нормоконтроля & 2 & 2 & 2 & 1 \\
\hline
Автоматическое обнаружение геометрических ошибок & 2 & 2 & 2 & 1 \\
\hline
\end{longtable}
}

\clearpage